\chapter{并发编程}

本章对应原书第~12~章 Concurrent Programming。课堂尚未讲授，以下仅保留目录。

\section{基于进程的并发编程}

\pending

\subsection{基于进程的并发服务器}

\pending

\subsection{进程的优劣}

\pending

\section{基于I/O多路复用的并发编程}

\pending

\subsection{基于I/O多路复用的并发事件驱动服务器}

\pending

\subsection{I/O多路复用的优劣}

\pending

\section{基于线程的并发编程}

\pending

\subsection{线程执行模型}

\pending

\subsection{Posix线程}

\pending

\subsection{创建线程}

\pending

\subsection{终止线程}

\pending

\subsection{回收已终止的线程}

\pending

\subsection{分离线程}

\pending

\subsection{初始化线程}

\pending

\subsection{基于线程的并发服务器}

\pending

\section{多线程程序中的共享变量}

\pending

\subsection{线程内存模型}

\pending

\subsection{将变量映射到内存}

\pending

\subsection{共享变量}

\pending

\section{用信号量同步线程}

\pending

\subsection{进度图}

\pending

\subsection{信号量}

\pending

\subsection{使用信号量来实现互斥}

\pending

\subsection{利用信号量来调度共享资源}

\pending

\subsection{综合：基于预线程化的并发服务器}

\pending

\section{使用线程提高并行性}

\pending

\section{其他并发问题}

\pending

\subsection{线程安全}

\pending

\subsection{可重入性}

\pending

\subsection{在线程化的程序中使用已存在的库函数}

\pending

\subsection{竞争}

\pending

\subsection{死锁}

\pending

\section{小结}

\pending
